



\section{Performance and Advantages of \textsc{GRACE}}

We evaluate \textsc{GRACE} across three principal axes: adversarial robustness, latent geometric structure, and reference-aware preference fidelity. Our experiments span 17 open-source LLMs and 12 attack types—including jailbreaks, logic inversions, and prompt injections—demonstrating consistent performance gains over DPO and its variants.

\vspace{1mm}
\subsection*{Adversarial Robustness: Lowering the Floor of Vulnerability}

Across the consolidated adversarial suite—comprising our benchmark corpus, Anthropic’s jailbreak dataset~\cite{perez2022red}, and prompt perturbations from PromptBench~\cite{zhu2024promptbench}—\textsc{GRACE} consistently lowers Attack Success Rate (ASR) compared to baselines. On models such as LLaMA-3 (8B), DeepSeek (7B), and Mixtral (8x22B), we observe ASR reductions of up to \textbf{30\%} post-training, with no degradation in performance on clean prompts.

\vspace{1mm}
\subsection*{Latent Geometry: Structural Interpretability and Generalization}

Using metrics like the Adversarial Vulnerability Quality Index (AVQI) and Density-Based Separation (DBS), we show that \textsc{GRACE} produces disentangled clusters in the pooled latent space:
\begin{itemize}[leftmargin=1.5em]
    \item \textbf{Safe completions} form low-variance clusters, well-separated from adversarial behavior.
    \item \textbf{Unsafe and jailbreak completions} coalesce into a compact adversarial manifold, distinct from the safe subspace.
\end{itemize}




\begin{figure*}[htp!]
    \centering
    \includegraphics[width=\textwidth]{figures/LLM_Attack_Benchmark_Heatmap_PostGRACE_OpenOnly_400dpi_OneLine.pdf}
    \caption{
    \textbf{GRACE Mitigation Performance across Open-Source LLMs.}
    This heatmap presents the Attack Success Rate (ASR) of 17 open-source large language models across 12 adversarial attack types. We show each attack the pre- and post-GRACE ASR, visually highlighting post-GRACE rows by golden borders. Each cell in a post-GRACE row displays the updated ASR (rounded) and the percentage reduction relative to the original, using a two-line format. GRACE (Geometric Preference Optimization for Alignment Robustness) consistently reduces ASR across all model scales and architectures—including instruction-tuned and chat-optimized variants like LLaMA-2/3, Vicuna, Mistral, Gemma, and DeepSeek—without requiring task-specific tuning or retraining. Attacks such as Goal Hijacking, Prompt Extraction, and TAP demonstrate substantial mitigation, reflecting GRACE’s strength against structural and semantic adversarial behaviors. This benchmark confirms that GRACE can provide effective, generalizable safety alignment while maintaining model usability.
    }
    \label{fig:grace_benchmark}
\end{figure*}

These geometric outcomes support the hypothesis that adversarial robustness arises from latent-space structure—not surface-level alignment.

\vspace{1mm}
\subsection*{KL Relaxation and Reference Drift Mitigation}

Direct Preference Optimization (DPO) often over-regularizes toward a fixed reference policy \(\pi_{\mathrm{ref}}\), risking underperformance when \(\pi_{\mathrm{ref}}\) itself produces unsafe outputs. \textsc{GRACE} relaxes this constraint via a tunable scaling factor \(\alpha \in [0,1]\)~\cite{wu2024generalized, chen2023epsilon}, allowing the model to:
\begin{itemize}[leftmargin=1.5em]
    \item Escape faulty reference completions while preserving overall alignment.
    \item Learn safer behaviors even when \(\pi_{\mathrm{ref}}\) is compromised.
\end{itemize}

This reduces KL-induced overfitting and improves generalization to adversarial contexts.

\vspace{1mm}
\subsection*{Lightweight and Modular Design}

\textsc{GRACE} requires no additional decoders or classifier heads. It operates entirely over frozen LLM representations and introduces only a soft attention profile \(\alpha^{(l)}\) over internal layers. This design ensures:
\begin{itemize}[leftmargin=1.5em]
    \item \textbf{Parameter efficiency} with minimal memory overhead.
    \item \textbf{Model agnosticity}—easily adaptable to any pretrained LLM.
    \item \textbf{Deployment ease} when using pre-trained \(\alpha^{(l)}\) vectors.
\end{itemize}

\vspace{1mm}
\subsection*{Summary of Core Advantages}

\begin{itemize}[leftmargin=1.8em]
    \item \textbf{Adversarial Robustness:} Up to 30\% ASR reduction across challenging attacks.
    \item \textbf{Latent Interpretability:} Behavior types form separated, analyzable clusters.
    \item \textbf{KL-Resilient Preference Learning:} Learns to prefer safe responses even with imperfect reference policies.
    \item \textbf{Modular and Lightweight:} No new architecture required—only learnable attention over frozen LLM layers.
\end{itemize}

\vspace{1mm}
In summary, \textsc{GRACE} unifies the strengths of preference modeling with the inductive bias of latent geometry, offering a scalable path toward adversarially aligned, interpretable, and mechanistically grounded language models.








